% =============================================================================
% Title    : SA Time Stability Tools - README (LuaLaTeX edition)
% Filename : README_LuaLaTeX.tex
% Version  : 0.0.3  (Semantic Versioning 2.0.0; same as the package)
% Author   : William W. Wallace
% Company  : NRAO / GBO
% Email    : wwallace@nrao.edu
% Created  : 2026-10-11
% Modified : 2026-10-11
% Engine   : LuaLaTeX (LuaHBTeX), TeX Live 2023 or later
% Build    : latexmk -lualatex README_LuaLaTeX.tex
%            (or run lualatex twice: contents, bookmarks and page totals)
% Licence  : GPL-3.0-or-later (SPDX), as the package
% Status   : Internal draft
%
% DESCRIPTION
%   The package README.md typeset with LuaLaTeX. The body was converted from
%   README.md with pandoc 3.1.3 and a Lua filter (headings shifted so
%   "##" = \section, inline code with break points, data-gap notes in boxes);
%   the preamble is hand-written. Edit README.md first and regenerate, or
%   edit this file directly - both are plain LaTeX.
%
%   PDF bookmarks: section > subsection hierarchy (hyperref + bookmark),
%   numbered, opened to the first level; the PDF opens with the bookmark
%   pane shown. Every reference URL and internal link is live.
%
%   Fonts: DejaVu Serif / Sans / Sans Mono with TeX Gyre DejaVu Math (all in
%   TeX Live). If they are missing the file falls back to TeX Gyre Pagella
%   and Latin Modern, so it builds on any current TeX Live or MiKTeX.
%
% REVISION HISTORY
%   0.0.3  2026-10-11  First LaTeX edition (README 0.0.3).
% =============================================================================
\documentclass[10pt,a4paper]{article}

% ---- page geometry ----------------------------------------------------------
\usepackage[a4paper,top=22mm,bottom=24mm,left=20mm,right=20mm,
            headheight=14pt]{geometry}

% ---- fonts (LuaLaTeX / fontspec / unicode-math) -------------------------------
\usepackage{amsmath}
\usepackage{fontspec}
\usepackage{unicode-math}
\IfFontExistsTF{DejaVu Serif}{%
  \setmainfont{DejaVu Serif}[Scale=0.95]
  \setsansfont{DejaVu Sans}[Scale=0.95]
  \setmonofont{DejaVu Sans Mono}[Scale=0.86,Ligatures=TeXOff]
}{%
  \setmainfont{TeX Gyre Pagella}
  \setsansfont{TeX Gyre Heros}
  \setmonofont{Latin Modern Mono}[Ligatures=TeXOff]
}
\IfFontExistsTF{TeX Gyre DejaVu Math}{\setmathfont{TeX Gyre DejaVu Math}[Scale=0.95]}{%
  \IfFontExistsTF{DejaVu Math TeX Gyre}{\setmathfont{DejaVu Math TeX Gyre}[Scale=0.95]}{%
    \setmathfont{Latin Modern Math}}}

% ---- colours, lists, tables, code ----------------------------------------------
\usepackage[dvipsnames,svgnames]{xcolor}
\definecolor{accent}{HTML}{2C5282}
\definecolor{accentlight}{HTML}{EBF4FF}
\definecolor{codebg}{HTML}{F4F6F8}
\definecolor{notebg}{HTML}{FFF8E6}
\definecolor{noteframe}{HTML}{C05621}
\usepackage{enumitem}
\setlist{itemsep=1pt,topsep=3pt}
\providecommand{\tightlist}{\setlength{\itemsep}{0pt}\setlength{\parskip}{0pt}}
\usepackage{longtable,booktabs,array,calc}
\usepackage{etoolbox}
\makeatletter
\patchcmd\longtable{\par}{\if@noskipsec\mbox{}\fi\par}{}{}
\makeatother
\IfFileExists{footnotehyper.sty}{\usepackage{footnotehyper}}{\usepackage{footnote}}
\makesavenoteenv{longtable}
% Table cells: ragged right with FINITE stretch (ragged2e), so a line is
% only broken early when it must - a plain \raggedright cell happily breaks
% a URL right after "https:".
\usepackage{ragged2e}
\AtBeginEnvironment{longtable}{\footnotesize\renewcommand{\arraystretch}{1.18}%
  \let\raggedright\RaggedRight}
\setlength{\LTpre}{6pt}
\setlength{\LTpost}{6pt}
\usepackage{fvextra}
\DefineVerbatimEnvironment{codeblock}{Verbatim}{%
  breaklines,breakanywhere,fontsize=\footnotesize,
  frame=leftline,framerule=2.5pt,rulecolor=\color{accent},
  framesep=6pt,xleftmargin=8pt}
\usepackage[most]{tcolorbox}
\newtcolorbox{notebox}{breakable,enhanced,colback=notebg,
  colframe=noteframe,boxrule=0pt,leftrule=3pt,arc=0pt,
  left=6pt,right=6pt,top=4pt,bottom=4pt,fontupper=\small}

% ---- headings, header/footer -----------------------------------------------------
\usepackage{needspace}  % keeps a heading with the start of its section
\usepackage{titlesec}
\titleformat{\section}{\sffamily\Large\bfseries\color{accent}}{\thesection}{0.7em}{}
  [{\color{accent}\titlerule[0.8pt]}]
\titleformat{\subsection}{\sffamily\large\bfseries\color{accent}}{\thesubsection}{0.6em}{}
\titleformat{\subsubsection}{\sffamily\normalsize\bfseries\color{accent}}{\thesubsubsection}{0.5em}{}
\titlespacing*{\section}{0pt}{14pt plus 4pt minus 2pt}{6pt}
\titlespacing*{\subsection}{0pt}{10pt plus 3pt minus 2pt}{4pt}
\setcounter{secnumdepth}{3}
\setcounter{tocdepth}{2}
\usepackage[titles]{tocloft}   % compact contents; number boxes fit 12.10
\setlength{\cftbeforesecskip}{3pt}
\setlength{\cftsecnumwidth}{2.0em}
\setlength{\cftsubsecindent}{2.0em}
\setlength{\cftsubsecnumwidth}{2.9em}
\renewcommand{\cftsecfont}{\sffamily\bfseries}
\renewcommand{\cftsecpagefont}{\sffamily\bfseries}
\usepackage{lastpage}
\usepackage{fancyhdr}
\pagestyle{fancy}
\fancyhf{}
\renewcommand{\headrulewidth}{0.4pt}
\fancyhead[L]{\sffamily\footnotesize\color{accent}SA Time Stability Tools 0.0.3}
\fancyhead[R]{\sffamily\footnotesize\nouppercase{\leftmark}}
\fancyfoot[C]{\sffamily\scriptsize SA Time Stability Tools 0.0.3 -- README (LuaLaTeX) --
  GPL-3.0-or-later \quad page \thepage\ of \pageref*{LastPage}}
\fancypagestyle{plain}{\fancyhf{}\renewcommand{\headrulewidth}{0pt}%
  \fancyfoot[C]{\sffamily\scriptsize SA Time Stability Tools 0.0.3 -- README (LuaLaTeX) --
  GPL-3.0-or-later \quad page \thepage\ of \pageref*{LastPage}}}

\setlength{\parindent}{0pt}
\setlength{\parskip}{4pt plus 1pt}
\setlength{\emergencystretch}{3em}
\usepackage{microtype}

% ---- hyperlinks and the PDF bookmark hierarchy ------------------------------------
\PassOptionsToPackage{hyphens}{url}
\usepackage[unicode]{hyperref}
\usepackage{bookmark}
\usepackage{xurl}
\urlstyle{same}
\hypersetup{
  pdftitle={SA Time Stability Tools 0.0.3 - Time Stability App - README},
  pdfauthor={William W. Wallace (NRAO/GBO)},
  pdfsubject={Frequency stability and phase-noise toolkit: user guide, equations, validation, references},
  pdfkeywords={Allan deviation, Theo1, ThêoH, NIST SP 1065, IEEE 1139, phase noise, Veusz, Stable32},
  pdfcreator={LuaLaTeX; body converted from README.md with pandoc 3.1.3},
  pdflang={en-US},
  pdfpagemode=UseOutlines,
  colorlinks=true, linkcolor=accent, urlcolor=accent, citecolor=accent,
}
\bookmarksetup{numbered, open, openlevel=1}

\begin{document}

% ---- title block ---------------------------------------------------------------
\thispagestyle{plain}
\pdfbookmark[1]{Title}{title}
\begin{center}
  {\sffamily\bfseries\huge\color{accent} SA Time Stability Tools 0.0.3\par}
  \vspace{4pt}
  {\sffamily\Large Time Stability App -- README\par}
  \vspace{8pt}
  {\normalsize William W. Wallace \quad NRAO / GBO \quad
   \href{mailto:wwallace@nrao.edu}{wwallace@nrao.edu}\par}
  \vspace{2pt}
  {\small Version 0.0.3 \quad 2026-10-11 \quad GPL-3.0-or-later \quad Internal draft\par}
\end{center}
\vspace{4pt}
{\color{accent}\hrule height 1.2pt}
\vspace{6pt}

\renewcommand{\contentsname}{Table of contents}
\pdfbookmark[1]{Table of contents}{contents}
{\small\tableofcontents}
\clearpage

% ---- body (generated from README.md) ----------------------------------------------
\Needspace{8\baselineskip}

\hypertarget{overview}{%
\section*{Overview}\label{overview}}
\addcontentsline{toc}{section}{Overview}

Python toolkit and desktop GUI for long-time-measurement frequency stability
and phase noise. It computes NIST SP 1065 time- and frequency-domain
statistics {[}1{]}, derives SSB phase noise, identifies power-law noise types and
frequency drift, converts vendor spectrum-analyzer exports, and writes
Stable32 files and native Veusz HDF5 (\texttt{.\allowbreak{}vszh5}) documents.

\begin{longtable}[]{@{}
  >{\raggedright\arraybackslash}p{(\columnwidth - 2\tabcolsep) * \real{0.1335}}
  >{\raggedright\arraybackslash}p{(\columnwidth - 2\tabcolsep) * \real{0.8665}}@{}}
\toprule\noalign{}
\endhead
\bottomrule\noalign{}
\endlastfoot
Version & \textbf{0.0.3} (Semantic Versioning 2.0.0 {[}27{]}; \texttt{0.0.x} = internal draft) \\
Date & 2026-10-11 \\
Author & William W. Wallace · NRAO/GBO · \href{mailto:wwallace@nrao.edu}{\nolinkurl{wwallace@nrao.edu}} \\
Licence & GPL-3.0-or-later (SPDX); full text in \texttt{LICENSE} \\
GUI & \texttt{Time\_\allowbreak{}Stability\_\allowbreak{}App.py} (Qt.py → PySide6) \\
Python & 3.8 -- 3.13. Suite run on 3.8.20, 3.12.13 and 3.13.16 \\
Veusz & Not needed to write the \texttt{.\allowbreak{}vszh5}; checked against Veusz 3.6.2 and 4.2.1 \\
Status & Internal draft, validated on synthetic data only \\
\end{longtable}

\textbf{0.0.3 is pass 2 of 2.} Pass 1 (0.0.2) moved the Veusz output to the native
HDF5 writer and renamed the GUI. Pass 2 fixes everything found in review,
parallelises the analysis and re-checks every reference. The full history is
in \texttt{CHANGELOG.md}.

\Needspace{8\baselineskip}

\hypertarget{what-changed-in-003}{%
\section{What changed in 0.0.3}\label{what-changed-in-003}}

\begin{longtable}[]{@{}
  >{\raggedright\arraybackslash}p{(\columnwidth - 6\tabcolsep) * \real{0.0420}}
  >{\raggedright\arraybackslash}p{(\columnwidth - 6\tabcolsep) * \real{0.4463}}
  >{\raggedright\arraybackslash}p{(\columnwidth - 6\tabcolsep) * \real{0.3835}}
  >{\raggedright\arraybackslash}p{(\columnwidth - 6\tabcolsep) * \real{0.1282}}@{}}
\toprule\noalign{}
\begin{minipage}[b]{\linewidth}\raggedright
\#
\end{minipage} & \begin{minipage}[b]{\linewidth}\raggedright
Problem (0.0.1/0.0.2)
\end{minipage} & \begin{minipage}[b]{\linewidth}\raggedright
Fix in 0.0.3
\end{minipage} & \begin{minipage}[b]{\linewidth}\raggedright
Checked by
\end{minipage} \\
\midrule\noalign{}
\endhead
\bottomrule\noalign{}
\endlastfoot
1 & A \texttt{freq\_\allowbreak{}hz}/\texttt{frac\_\allowbreak{}freq} CSV loaded with \texttt{len(x) =\allowbreak{} len(t)+1}; drift removal, plots and export failed & Time tags extended by one interval; \texttt{load\_\allowbreak{}canonical\_\allowbreak{}series()} reports the column read & \texttt{validate\_\allowbreak{}core} T9, \texttt{validate\_\allowbreak{}gui} \\
2 & Loading a new file kept the old drift fit and results (stale residual applied) & Results cleared on load; late results of a running job discarded & \texttt{validate\_\allowbreak{}gui} \\
3 & GUI always logged "converting to phase"; README said the loader returned the column type & Log names the column; f0 filled from \texttt{freq\_\allowbreak{}hz} data & \texttt{validate\_\allowbreak{}gui} \\
4 & Decimal-comma exports (\texttt{;}, tab) read as "no numeric data" & Delimiter chosen from the data lines; decimal commas read & \texttt{validate\_\allowbreak{}sa\_\allowbreak{}convert} \\
5 & Converted CSV kept 11 significant digits (10 GHz series: 3.7 mHz error on 1.2 mHz steps) & Shortest exact decimal (bit-exact round trip) & \texttt{validate\_\allowbreak{}sa\_\allowbreak{}convert} \\
6 & Stable32 writer left a partial file when it rejected a last-point gap & Validation before writing; atomic \texttt{.part} + rename & \texttt{validate\_\allowbreak{}stable32\_\allowbreak{}export} \\
7 & \texttt{NOISE\_\allowbreak{}TABLE} "Linear Freq. Drift" +1.5 was Random-Run FM; drift (+1) reported as Flicker-Walk FM & Row renamed; drift detected from OADEV +1 with a flat/falling OHDEV & \texttt{validate\_\allowbreak{}core} T10 \\
8 & Drift rate logged as "s/s\^{}2" & \(D = 2a_2\) in 1/s & \texttt{validate\_\allowbreak{}core} T11 \\
9 & Everything single-threaded on the GUI thread; ThêoH of 10⁵ points took 218 s & Worker threads in the GUI; faster Thêo1 kernel; thread/process pools; optional CuPy & \texttt{validate\_\allowbreak{}core} T8, below \\
10 & References garbled (wrong DOI, pages, authors, merged citations) & All checked against primary records; table below & --- \\
11 & Three validation scripts always exited 0 & Every script returns PASS/FAIL; \texttt{run\_\allowbreak{}all.py} & \texttt{run\_\allowbreak{}all.py} \\
12 & \texttt{PySide6\textasciitilde{}=\allowbreak{}6.6.3} pin: no Python 3.13 wheels & Version-gated pins & install \\
13 & GUI: no Qt.py, no menus, no light/dark display & Qt.py; File/Analysis/Export/View/Help; View › Light/Dark (remembered) & \texttt{validate\_\allowbreak{}gui} \\
--- & Phase-noise ID from two-bin differences named the right type in 67 \% of segments & Windowed least-squares slope: 99 \% & \texttt{validate\_\allowbreak{}phase\_\allowbreak{}noise} \\
--- & ThêoH stopped at 62.8 \% of the record & Grid ends at the largest even \(m \le N-1\): 75 \% & \texttt{validate\_\allowbreak{}core} T6 \\
--- & Long records drawn point by point & Min/max envelope for display (GUI and Veusz pages); every sample analysed and stored & \texttt{validate\_\allowbreak{}veusz\_\allowbreak{}export} \\
\end{longtable}

\Needspace{8\baselineskip}

\hypertarget{contents}{%
\section{Contents}\label{contents}}

\begin{longtable}[]{@{}
  >{\raggedright\arraybackslash}p{(\columnwidth - 4\tabcolsep) * \real{0.2045}}
  >{\raggedright\arraybackslash}p{(\columnwidth - 4\tabcolsep) * \real{0.0970}}
  >{\raggedright\arraybackslash}p{(\columnwidth - 4\tabcolsep) * \real{0.6985}}@{}}
\toprule\noalign{}
\begin{minipage}[b]{\linewidth}\raggedright
File
\end{minipage} & \begin{minipage}[b]{\linewidth}\raggedright
Version
\end{minipage} & \begin{minipage}[b]{\linewidth}\raggedright
Purpose
\end{minipage} \\
\midrule\noalign{}
\endhead
\bottomrule\noalign{}
\endlastfoot
\texttt{Time\_\allowbreak{}Stability\_\allowbreak{}App.py} & 0.0.3 & Qt GUI (Qt.py → PySide6) \\
\texttt{stability\_\allowbreak{}core.py} & 0.0.3 & CSV I/O, drift fit, ADEV/OADEV/MDEV/TDEV/OHDEV/Thêo1/ThêoBR/ThêoH, noise and drift identification, thread/process/GPU engine \\
\texttt{phase\_\allowbreak{}noise.py} & 0.0.3 & SSB L(f) via Welch PSD, direct traces, frequency-domain noise ID, L(f)↔S\_y(f) \\
\texttt{sa\_\allowbreak{}convert.py} & 0.0.3 & Keysight / R\&S / generic exports → canonical CSV; parallel batch marker extraction \\
\texttt{stable32\_\allowbreak{}export.py} & 0.0.3 & Stable32 phase, frequency and sigma-tau files \\
\texttt{veusz\_\allowbreak{}export.py} & 0.0.3 & Native Veusz HDF5 export, display envelope, optional render \\
\texttt{vszh5\_\allowbreak{}writer.py} & 0.0.5 & Shared Veusz HDF5 writer, unchanged and byte-identical to QuickSyn 0.0.5 (SHA-256 \texttt{2cdf43d2…ae33c1c}); see \protect\hyperlink{notes-for-the-shared-writer-vszh5_writerpy-005}{Notes for the shared writer} \\
\texttt{validate\_\allowbreak{}core.py} & 0.0.3 & 11 tests: estimators vs allantools, slopes, Thêo1, ThêoH grid, parallel = serial, CSV loader, drift \\
\texttt{validate\_\allowbreak{}phase\_\allowbreak{}noise.py} & 0.0.3 & L(f) slopes (independent seeds), noise-ID accuracy, S\_y(f) \\
\texttt{validate\_\allowbreak{}sa\_\allowbreak{}convert.py} & 0.0.3 & Vendor files incl. decimal commas, precision, parallel batch \\
\texttt{validate\_\allowbreak{}stable32\_\allowbreak{}export.py} & 0.0.3 & New: format, gaps, atomic write, sigma-tau table \\
\texttt{validate\_\allowbreak{}veusz\_\allowbreak{}export.py} & 0.0.3 & \texttt{.\allowbreak{}vszh5} layout, bit-exact data, envelope --- no Veusz needed \\
\texttt{validate\_\allowbreak{}vszh5\_\allowbreak{}in\_\allowbreak{}veusz.py} & 0.0.3 & Round trip through Veusz\textquotesingle s own loader and renderer (needs Veusz) \\
\texttt{validate\_\allowbreak{}gui.py} & 0.0.3 & Headless GUI test: menus, themes, worker threads, exports \\
\texttt{run\_\allowbreak{}all.py} & 0.0.3 & New: runs every \texttt{validate\_\allowbreak{}*.py} in parallel, summary table, exit status \\
\texttt{environment.yml} / \texttt{requirements.txt} & 0.0.3 & mamba (conda-forge) environment / pip alternative \\
\texttt{README.md}, \texttt{CHANGELOG.md}, \texttt{LICENSE} & 0.0.3 & Documentation (Markdown with \texttt{\$…\$} LaTeX for Obsidian); GPL-3.0 text \\
\texttt{README.pdf} & 0.0.3 & This README printed from HTML (Chromium); PDF bookmarks follow the headings \\
\texttt{README\_\allowbreak{}LuaLaTeX.pdf}, \texttt{latex/\allowbreak{}README\_\allowbreak{}LuaLaTeX.tex} & 0.0.3 & The README typeset with LuaLaTeX: numbered sections, contents, bookmark hierarchy, linked references. Rebuild with \texttt{latexmk -lualatex README\_\allowbreak{}LuaLaTeX.tex} \\
\texttt{example\_\allowbreak{}data/\allowbreak{}example\_\allowbreak{}maser\_\allowbreak{}phase.csv} & --- & Example record (8193 points, \texttt{time\_\allowbreak{}s,\allowbreak{}phase\_\allowbreak{}s}). Synthetic in scale (x of order 10³ s), so L(f) lies far above 0 dBc/Hz; for workflow checks only \\
\end{longtable}

\Needspace{8\baselineskip}

\hypertarget{installation-mamba-conda-forge}{%
\section{Installation (mamba, conda-forge)}\label{installation-mamba-conda-forge}}

\begin{codeblock}
mamba env create -f environment.yml
mamba activate sa_stability
python Time_Stability_App.py
python run_all.py --veusz-python python      # optional: the validation suite
\end{codeblock}

\texttt{h5py} writes the \texttt{.\allowbreak{}vszh5}; Veusz is in the environment only to open the
documents and render PDF/PNG (remove it if you use the Veusz Windows
installer, then set \texttt{TIME\_\allowbreak{}STABILITY\_\allowbreak{}VEUSZ} to \texttt{veusz.exe}). \texttt{threadpoolctl}
lets the Thêo1 thread pool cap BLAS threads (without it a BLAS-free kernel is
used). For the GPU path add \texttt{cupy} (or \texttt{cupy-cuda12x} with pip) to match your
CUDA driver.

pip alternative: \texttt{python -m pip install -r requirements.txt}.

\Needspace{8\baselineskip}

\hypertarget{running-the-gui}{%
\section{Running the GUI}\label{running-the-gui}}

\begin{longtable}[]{@{}
  >{\raggedright\arraybackslash}p{(\columnwidth - 2\tabcolsep) * \real{0.1080}}
  >{\raggedright\arraybackslash}p{(\columnwidth - 2\tabcolsep) * \real{0.8920}}@{}}
\toprule\noalign{}
\begin{minipage}[b]{\linewidth}\raggedright
Menu
\end{minipage} & \begin{minipage}[b]{\linewidth}\raggedright
Items (shortcut)
\end{minipage} \\
\midrule\noalign{}
\endhead
\bottomrule\noalign{}
\endlastfoot
File & Load Canonical CSV (Ctrl+O), Convert Vendor SA File, Close Data --- frees memory (Ctrl+W), Quit (Ctrl+Q) \\
Analysis & Remove Drift (Ctrl+D), Run Stability Analysis (Ctrl+T), Run SSB Phase Noise (Ctrl+P), Run All (Ctrl+R) \\
Export & Stable32 file --- phase \texttt{.PHD}, frequency \texttt{.FRD} or sigma-tau table \texttt{.DAT}; Veusz HDF5 \texttt{.\allowbreak{}vszh5} (Ctrl+E); Render a \texttt{.\allowbreak{}vszh5} to PDF/PNG \\
View & Light display / Dark display (Fusion palette and plot colours; remembered) \\
Help & About (versions, Qt binding, CPUs, GPU) \\
\end{longtable}

Workflow: load a canonical CSV (or convert an analyzer export and say whether
the series holds \texttt{phase\_\allowbreak{}s}, \texttt{freq\_\allowbreak{}hz} or \texttt{frac\_\allowbreak{}freq}) → check \texttt{tau0} (from the
time stamps) and \texttt{f0} (filled from \texttt{freq\_\allowbreak{}hz} data) → Run All, or drift →
stability → phase noise → export. Analyses run on a worker thread: the window
stays live, buttons are disabled while a job runs, and the status bar shows
the job, the thread count and the GPU state. \textbf{Threads} sets the workers of
the analysis engine (default: all CPUs); \textbf{GPU (CuPy)} is enabled when CuPy
is installed.

\Needspace{5\baselineskip}

\hypertarget{spyder}{%
\subsection{Spyder}\label{spyder}}

Every file is split into \texttt{\# \%\%} cells (header, imports, one cell per section,
entry point, a commented scratch cell) and runs cell by cell in order ---
checked with a harness that reproduces Spyder\textquotesingle s \texttt{runcell} namespace rules
(\texttt{\_\allowbreak{}\_\allowbreak{}file\_\allowbreak{}\_} defined only while a cell runs). Scripts capture their path once,
never call \texttt{sys.\allowbreak{}exit()} inside IPython, and the GUI reuses an existing
\texttt{QApplication}. With the console\textquotesingle s graphics backend set to Qt the window stays
live and \texttt{main()} returns at once; with the inline backend \texttt{main()} runs the
event loop until the window closes. The process pools never re-run your
script or console cell in their workers (see Performance).

\Needspace{8\baselineskip}

\hypertarget{canonical-csv-format}{%
\section{Canonical CSV format}\label{canonical-csv-format}}

\begin{codeblock}
time_s,phase_s        # phase / time-error data (seconds), OR
time_s,freq_hz        # absolute frequency readings (Hz), OR
time_s,frac_freq      # fractional frequency y = (f - f0)/f0
\end{codeblock}

\texttt{stability\_\allowbreak{}core.\allowbreak{}load\_\allowbreak{}canonical\_\allowbreak{}series()} reads the file with NumPy\textquotesingle s fast
parser (a UTF-8 byte-order mark, \texttt{\#} comments and any column order are
accepted), sorts by time and returns the phase, the time tags (same length),
the column read, \texttt{tau0} (median spacing) and the \texttt{f0} used. Frequency
columns are integrated to phase (SP 1065 §4.1): N readings give N + 1 phase
points, and reading k is taken to cover \([t_k, t_k+\tau_0]\). \texttt{freq\_\allowbreak{}hz} is
normalised by its median unless \texttt{f0\_\allowbreak{}hz} is given; a constant fractional
offset changes no AVAR-family statistic. \texttt{load\_\allowbreak{}canonical\_\allowbreak{}csv()} keeps its
0.0.1 return form \texttt{(t,\allowbreak{} x,\allowbreak{} 'phase')}.

\Needspace{8\baselineskip}

\hypertarget{vendor-spectrum-analyzer-conversion--important-caveat}{%
\section{Vendor spectrum-analyzer conversion --- important caveat}\label{vendor-spectrum-analyzer-conversion--important-caveat}}

\textbf{Exact SA export layouts vary by instrument model and firmware.} Export one
small file from your instrument and confirm the units before trusting a bulk
conversion. \texttt{sa\_\allowbreak{}convert.py} provides:

\begin{itemize}
\tightlist
\item
  a generic parser for comma, semicolon, tab or white-space files with
  metadata/comment lines; the delimiter is chosen from the data lines, and
  with \texttt{;}, tab or white space a decimal comma (\texttt{1,\allowbreak{}234}) is read as 1.234
  (with a comma delimiter a decimal comma is ambiguous and not supported);
\item
  Keysight- and R\&S-flavoured parsers that take the numeric block after the
  instrument\textquotesingle s header (trailing \texttt{;} allowed); header words such as
  "phase noise" or "dBc" mark a phase-noise trace;
\item
  time-series vs trace classification of column 0 (uniform vs log-uniform
  spacing);
\item
  output at full double precision (shortest exact decimal); the value column
  of a time series must be named by the caller (\texttt{phase\_\allowbreak{}s}, \texttt{freq\_\allowbreak{}hz},
  \texttt{frac\_\allowbreak{}freq}) because the file does not say what the analyzer logged;
\item
  \texttt{batch\_\allowbreak{}extract\_\allowbreak{}marker\_\allowbreak{}time\_\allowbreak{}series()} for repeated snapshots (peak or
  n-th largest value per file), parsed in a process pool for 8 or more files.
\end{itemize}

\Needspace{8\baselineskip}

\hypertarget{stable32-export}{%
\section{Stable32 export}\label{stable32-export}}

Tab-delimited ASCII: \texttt{timetag  value} (or the value column only with
\texttt{timetags=\allowbreak{}False}), 16 significant digits, "gap = 0" with a gap count in the
\texttt{\#} header, no gap allowed at the first or last point --- checked \textbf{before}
anything is written, and the file is written to \texttt{<name>.part} and renamed, so
an error never leaves a truncated file {[}14{]}. The GUI exports the
(drift-removed) phase \texttt{.PHD}, its fractional frequency \texttt{.FRD} or the
sigma-tau table \texttt{.DAT} (τ, m, n, ADEV, OADEV, MDEV, TDEV, OHDEV; rows with a
non-finite value are skipped and counted).

\begin{notebox}

\textbf{Data gap --- time tags.} The time-tag column is elapsed seconds (as in
0.0.1). Whether Stable32 reads numeric time tags as MJD days when it derives
τ from them could not be confirmed from the documentation available here.
Enter τ in Stable32 yourself, or export without time tags. Reading the files
in Stable32 itself has not been tested.

\end{notebox}

\Needspace{8\baselineskip}

\hypertarget{veusz-export-vszh5}{%
\section{\texorpdfstring{Veusz export (\texttt{.vszh5})}{Veusz export (.vszh5)}}\label{veusz-export-vszh5}}

\texttt{veusz\_\allowbreak{}export.\allowbreak{}export\_\allowbreak{}vszh5()} writes one native Veusz HDF5 document with h5py
through \texttt{vszh5\_\allowbreak{}writer.py} {[}17{]} --- no Veusz and no Qt are imported, so the GUI
writes it in-process. Every result is stored at float64; the write is atomic.

\begin{longtable}[]{@{}
  >{\raggedright\arraybackslash}p{(\columnwidth - 4\tabcolsep) * \real{0.4383}}
  >{\raggedright\arraybackslash}p{(\columnwidth - 4\tabcolsep) * \real{0.4427}}
  >{\raggedright\arraybackslash}p{(\columnwidth - 4\tabcolsep) * \real{0.1190}}@{}}
\toprule\noalign{}
\begin{minipage}[b]{\linewidth}\raggedright
Dataset
\end{minipage} & \begin{minipage}[b]{\linewidth}\raggedright
Content
\end{minipage} & \begin{minipage}[b]{\linewidth}\raggedright
Tag
\end{minipage} \\
\midrule\noalign{}
\endhead
\bottomrule\noalign{}
\endlastfoot
\texttt{raw/\allowbreak{}time\_\allowbreak{}s}, \texttt{raw/\allowbreak{}phase\_\allowbreak{}s}, \texttt{raw/\allowbreak{}residual\_\allowbreak{}s} & Input and drift-removed x(t), every sample & raw \\
\texttt{raw/\allowbreak{}freq\_\allowbreak{}time\_\allowbreak{}s}, \texttt{raw/\allowbreak{}frac\_\allowbreak{}freq}, \texttt{raw/\allowbreak{}frac\_\allowbreak{}freq\_\allowbreak{}residual} & \(y_k = (x_{k+1}-x_k)/\tau_0\) at interval midpoints & raw \\
\texttt{plot/\allowbreak{}<series>\_\allowbreak{}time\_\allowbreak{}s}, \texttt{plot/\allowbreak{}<series>} & Min/max display envelope when a series is longer than 200 000 points & plot \\
\texttt{stability/\allowbreak{}tau\_\allowbreak{}s,\allowbreak{} m,\allowbreak{} n\_\allowbreak{}used,\allowbreak{} adev,\allowbreak{} oadev,\allowbreak{} mdev,\allowbreak{} ohdev,\allowbreak{} tdev\_\allowbreak{}s} & AVAR-family sweep & stability \\
\texttt{stability/\allowbreak{}theoh\_\allowbreak{}tau\_\allowbreak{}s,\allowbreak{} theoh\_\allowbreak{}dev,\allowbreak{} theoh\_\allowbreak{}m,\allowbreak{} theoh\_\allowbreak{}source} & ThêoH curve and the estimator of each point & stability \\
\texttt{noise\_\allowbreak{}id/\allowbreak{}*} (incl. \texttt{slope\_\allowbreak{}hdev}), \texttt{phase\_\allowbreak{}noise\_\allowbreak{}id/\allowbreak{}*} & Noise-identification tables & noise\_id \\
\texttt{phase\_\allowbreak{}noise/\allowbreak{}freq\_\allowbreak{}hz}, \texttt{phase\_\allowbreak{}noise/\allowbreak{}l\_\allowbreak{}dbc\_\allowbreak{}hz} & SSB L(f) & phase\_noise \\
\texttt{curves/\allowbreak{}*} & Extra curves passed as \texttt{sigma\_\allowbreak{}tau\_\allowbreak{}series} (0.0.1 API) & curves \\
\texttt{meta/\allowbreak{}title,\allowbreak{} source,\allowbreak{} generator,\allowbreak{} parameters\_\allowbreak{}json,\allowbreak{} datasets} & Run description, tau0, f0, drift fit, envelope flag & meta \\
\end{longtable}

Pages: \texttt{page\_\allowbreak{}overview} (2 × 2), \texttt{page\_\allowbreak{}phase}, \texttt{page\_\allowbreak{}frequency},
\texttt{page\_\allowbreak{}sigma\_\allowbreak{}tau} (ADEV, OADEV, MDEV, OHDEV, ThêoH), \texttt{page\_\allowbreak{}tdev} (seconds, on
its own axes), \texttt{page\_\allowbreak{}phase\_\allowbreak{}noise}. Long series are drawn from the envelope
(each bin keeps its minimum and maximum, in time order, so no excursion is
lost); Veusz stays responsive with a few hundred thousand points per trace.

\Needspace{5\baselineskip}

\hypertarget{rendering-optional}{%
\subsection{Rendering (optional)}\label{rendering-optional}}

\texttt{render\_\allowbreak{}vszh5()} / \texttt{python veusz\_\allowbreak{}export.\allowbreak{}py --render FILE.\allowbreak{}vszh5} writes one
PDF of all pages and one PNG per page. The Veusz program is driven in
\texttt{--listen} mode with its HDF5 loader (\texttt{veusz --export} cannot open a
\texttt{.\allowbreak{}vszh5}), with the veusz Python package as fallback; the program is searched
in \texttt{TIME\_\allowbreak{}STABILITY\_\allowbreak{}VEUSZ}, \texttt{VEUSZ\_\allowbreak{}EXE}, \texttt{PATH}. The GUI renders in a child
process, so Veusz\textquotesingle s Qt never loads next to PySide6.

\Needspace{8\baselineskip}

\hypertarget{performance-and-parallelism}{%
\section{Performance and parallelism}\label{performance-and-parallelism}}

\begin{longtable}[]{@{}
  >{\raggedright\arraybackslash}p{(\columnwidth - 2\tabcolsep) * \real{0.1300}}
  >{\raggedright\arraybackslash}p{(\columnwidth - 2\tabcolsep) * \real{0.8700}}@{}}
\toprule\noalign{}
\begin{minipage}[b]{\linewidth}\raggedright
Stage
\end{minipage} & \begin{minipage}[b]{\linewidth}\raggedright
How it runs
\end{minipage} \\
\midrule\noalign{}
\endhead
\bottomrule\noalign{}
\endlastfoot
GUI & Analyses on a worker thread; results delivered by a Qt signal; stale results dropped \\
Thêo1 kernel & Array views and two reused buffers per averaging factor --- no temporaries, no index arrays \\
ThêoH & Every Thêo1 the curve needs (bias ladder + grid) in one batch, largest first: a spawn \textbf{process pool} for ≥ 100 000 points, a \textbf{thread pool} for ≥ 50 000, serial below (measured: threads lose there) \\
One large Thêo1 & Split over the w-range on threads \\
BLAS & Capped to one thread inside the pools (threadpoolctl). Uncapped, OpenBLAS\textquotesingle s own threads oversubscribe the CPUs: one Thêo1 point (N = 100 000, m = 20 000) took 3.89 s on 2 threads vs 1.48 s serial and 0.95 s on 2 threads with the cap \\
AVAR sweep & One thread per averaging factor for long records \\
GPU & \texttt{use\_\allowbreak{}gpu=\allowbreak{}True} / GUI \textbf{GPU (CuPy)}: Thêo1 sums with CuPy \\
Welch FFT & \texttt{scipy.fft} worker threads \\
Batch conversion & Process pool (thread fallback) \\
Display & Min/max envelope (50 000 points on screen, 200 000 in Veusz) \\
\end{longtable}

Measured on the 2-vCPU test VM (whose ceiling for perfectly parallel NumPy
work is 1.47×), white/flicker FM, N = 100 000:

\begin{longtable}[]{@{}llll@{}}
\toprule\noalign{}
Quantity & 0.0.2 & 0.0.3 serial & 0.0.3 parallel \\
\midrule\noalign{}
\endhead
\bottomrule\noalign{}
\endlastfoot
Thêo1, m = 2 000 & 1.26 s & 0.23 s & 0.15 s (2 threads) \\
Thêo1, m = 20 000 & 8.31 s & 1.75 s & 1.07 s (2 threads) \\
Thêo1, m = 50 000 & 11.84 s & 2.64 s & 1.79 s (2 threads) \\
Full ThêoH curve & 218.1 s & 26.5 s & 23.0 s (2 processes) \\
\end{longtable}

Results are identical to the serial path to rounding (≤ 2 × 10⁻¹⁵) and to
0.0.2 (Thêo1 agrees with \texttt{allantools.\allowbreak{}theo1} to ≤ 6.6 × 10⁻¹⁶ on the benchmark
records, 2.5 × 10⁻¹⁶ in \texttt{validate\_\allowbreak{}core}). Larger gains
need more cores; they were not measurable here. The GPU path is written
against CuPy\textquotesingle s NumPy API and has \textbf{not} been run (no GPU here).

Spawned workers normally re-execute the parent\textquotesingle s main script; while a pool
starts, \texttt{\_\allowbreak{}\_\allowbreak{}main\_\allowbreak{}\_}\textquotesingle s \texttt{\_\allowbreak{}\_\allowbreak{}file\_\allowbreak{}\_}/\texttt{\_\allowbreak{}\_\allowbreak{}spec\_\allowbreak{}\_} are hidden from multiprocessing,
so a cell-style script, a Spyder console or the GUI is not re-run (or
re-imported with Qt) in each worker {[}28{]}.

\Needspace{8\baselineskip}

\hypertarget{time-domain-noise-identification}{%
\section{Time-domain noise identification}\label{time-domain-noise-identification}}

\(\sigma_y^2(\tau) \propto \tau^{\mu}\) with \(S_y(f) \propto f^{\alpha}\) and
\(\mu = -\alpha - 1\) for \(-4 \le \alpha < 1\) (\(\mu \approx -2\) for
\(\alpha \ge 1\)); the ADEV log-log slope is \(\mu/2\) {[}1{]}, {[}3{]}:

\begin{longtable}[]{@{}llll@{}}
\toprule\noalign{}
Noise type & α & ADEV slope & MDEV slope \\
\midrule\noalign{}
\endhead
\bottomrule\noalign{}
\endlastfoot
White PM & 2 & −1.0 & −1.5 \\
Flicker PM & 1 & −1.0 & −1.0 \\
White FM & 0 & −0.5 & −0.5 \\
Flicker FM & −1 & 0.0 & 0.0 \\
Random Walk FM & −2 & +0.5 & +0.5 \\
Flicker Walk FM & −3 & +1.0 † & +1.0 † \\
Random Run FM & −4 & +1.5 † & +1.5 † \\
Linear frequency drift & --- & +1.0 & +1.0 \\
\end{longtable}

† AVAR and MVAR converge only for \(\alpha > -3\); these rows are the Hadamard
(HDEV) slopes, which converge for \(\alpha > -5\). 0.0.1/0.0.2 labelled the +1.5
row "Linear Freq. Drift". A drift \(D\) gives \(\sigma_y(\tau) = D\tau/\sqrt{2}\)
(slope +1, the same as Flicker-Walk FM), while the Hadamard deviation does not
see it: \texttt{classify\_\allowbreak{}noise()} reports \textbf{Linear Freq. Drift} when the OADEV slope
is within 0.3 of +1 and the OHDEV slope is at least 0.5 lower (or OHDEV is
zero). White PM vs Flicker PM is separated by the MDEV slope (−1.5 vs −1.0).

\Needspace{8\baselineskip}

\hypertarget{frequency-domain-ssb-phase-noise-identification}{%
\section{Frequency-domain (SSB phase noise) identification}\label{frequency-domain-ssb-phase-noise-identification}}

\begin{longtable}[]{@{}ll@{}}
\toprule\noalign{}
Noise type & L(f) slope (dB/decade) \\
\midrule\noalign{}
\endhead
\bottomrule\noalign{}
\endlastfoot
White PM & 0 \\
Flicker PM & −10 \\
White FM & −20 \\
Flicker FM & −30 \\
Random Walk FM & −40 \\
\end{longtable}

\(\mathcal{L}(f) = 10\log_{10}[\tfrac12 S_\phi(f)]\) is already in dB, so its
slope against \(\log_{10} f\) is 10× the \(S_\phi\) exponent, with
\(S_\phi = (f_0/f)^2 S_y\) {[}3{]}. L(f) is averaged in log-spaced bins; from 0.0.3
each segment\textquotesingle s slope is a least-squares fit over ±0.25 decade of bins. On 15
synthetic records (5 noise types × 3 seeds, 0.05--20 Hz) the 0.0.2 two-bin
method named the injected type in 47--84 \% of segments (mean 67 \%), the
windowed fit in 95--100 \% (mean 99 \%). With independent seeds the measured
slopes are −19.80 / −29.88 / −40.06 dB/decade for white / flicker /
random-walk FM; the −0.8 dB/decade offset reported for 0.0.1 came from one
random realisation shared by all three cases.

\Needspace{8\baselineskip}

\hypertarget{equation-appendix}{%
\section{Equation appendix}\label{equation-appendix}}

Equations use \texttt{\$\$ … \$\$} and \texttt{\$ … \$} so they render in Obsidian. \(x_i\) are
phase samples at spacing \(\tau_0\), \(N\) their number, \(\tau = m\tau_0\).

\Needspace{5\baselineskip}

\hypertarget{overlapping-allan-deviation-oadev}{%
\subsection{Overlapping Allan deviation (OADEV)}\label{overlapping-allan-deviation-oadev}}

SP 1065 {[}1{]} Eq. (11):

\[
\sigma_y^2(m\tau_0) = \frac{1}{2(m\tau_0)^2 (N-2m)} \sum_{i=1}^{N-2m}
\left(x_{i+2m} - 2x_{i+m} + x_i\right)^2
\]

\Needspace{5\baselineskip}

\hypertarget{modified-allan-deviation-mdev}{%
\subsection{Modified Allan deviation (MDEV)}\label{modified-allan-deviation-mdev}}

SP 1065 {[}1{]} Eq. (14):

\[
\mathrm{Mod}\,\sigma_y^2(m\tau_0) = \frac{1}{2 m^2 (m\tau_0)^2 (N-3m+1)}
\sum_{j=1}^{N-3m+1}\left[\sum_{i=j}^{j+m-1}
\left(x_{i+2m}-2x_{i+m}+x_i\right)\right]^2
\]

\Needspace{5\baselineskip}

\hypertarget{time-deviation-tdev}{%
\subsection{Time deviation (TDEV)}\label{time-deviation-tdev}}

SP 1065 {[}1{]} Eq. (15):

\[
\sigma_x(\tau) = \frac{\tau}{\sqrt{3}}\,\mathrm{Mod}\,\sigma_y(\tau)
\]

\Needspace{5\baselineskip}

\hypertarget{overlapping-hadamard-deviation-ohdev}{%
\subsection{Overlapping Hadamard deviation (OHDEV)}\label{overlapping-hadamard-deviation-ohdev}}

SP 1065 {[}1{]} Eq. (20):

\[
\sigma_H^2(m\tau_0) = \frac{1}{6(m\tau_0)^2 (N-3m)}
\sum_{i=1}^{N-3m}\left(x_{i+3m}-3x_{i+2m}+3x_{i+m}-x_i\right)^2
\]

\Needspace{5\baselineskip}

\hypertarget{thuxeao1}{%
\subsection{Thêo1}\label{thuxeao1}}

SP 1065 {[}1{]} Eq. (30) and {[}5{]}; \(m\) even, \(10 \le m \le N-1\). SP 1065 sums
over \(\delta = 0 \dots m/2-1\); with \(w = m/2 - \delta\) and a 0-based \(i\) the
sum reads:

\[
\mathrm{Theo1}(m,\tau_0,N) = \frac{1}{0.75\,(N-m)\,(m\tau_0)^2}
\sum_{i=0}^{N-m-1}\sum_{w=1}^{m/2}
\frac{\left[(x_i - x_{i+w}) + (x_{i+m} - x_{i+m-w})\right]^2}{w}
\]

reported at the stride ("Theo1 averages every permissible squared
second-difference \ldots{} in a given span or stride \(\tau_s = 0.75\,m\,\tau_0\)" {[}5{]}):

\[
\tau = 0.75\, m\, \tau_0
\]

\Needspace{5\baselineskip}

\hypertarget{thuxeaobr-and-thuxeaoh}{%
\subsection{ThêoBR and ThêoH}\label{thuxeaobr-and-thuxeaoh}}

Howe {[}4{]} Eqs. (6)--(7):

\[
G = \frac{1}{n+1}\sum_{i=0}^{n}
\frac{\mathrm{AVAR}(m = 9+3i)}{\mathrm{Theo1}(m = 12+4i)},\qquad
n = \left\lfloor \frac{0.1N}{3} - 3 \right\rfloor,\qquad
\mathrm{Th\hat{e}oBR}(m) = \sqrt{G}\cdot\mathrm{Theo1}(m)
\]

\[
\mathrm{Th\hat{e}oH} =
\begin{cases}
\mathrm{AVAR}(m), & 1 \le m < k/\tau_0 \\
\mathrm{Th\hat{e}oBR}(m), & k/(0.75\,\tau_0) \le m \le N-1,\ m \text{ even}
\end{cases}
\]

with \(k\) the largest τ ≤ 10 \% of the record. The code subsamples the ladder
to at most 40 evenly spaced rungs and grows the ThêoBR grid by 1.3× per step,
ending at the largest even \(m \le N-1\) (75 \% of the record).

\Needspace{5\baselineskip}

\hypertarget{ssb-phase-noise}{%
\subsection{SSB phase noise}\label{ssb-phase-noise}}

IEEE Std 1139 {[}3{]}:

\[
\phi(t) = 2\pi f_0\, x(t), \qquad
\mathcal{L}(f) = 10\log_{10}\!\left[\tfrac{1}{2}\,S_\phi(f)\right]\ \mathrm{dBc/Hz}
\]

with \(S_\phi(f)\) the one-sided Welch {[}11{]} PSD of \(\phi(t)\), and {[}3{]}

\[
S_\phi(f) = 2\cdot 10^{\mathcal{L}(f)/10}, \qquad
S_y(f) = \frac{f^2}{f_0^2}\, S_\phi(f)
\]

\Needspace{5\baselineskip}

\hypertarget{fractional-frequency-from-phase-and-phase-from-frequency}{%
\subsection{Fractional frequency from phase, and phase from frequency}\label{fractional-frequency-from-phase-and-phase-from-frequency}}

SP 1065 {[}1{]} §4.1:

\[
y_k = \frac{x_{k+1} - x_k}{\tau_0}, \qquad
x_0 = 0,\quad x_{k+1} = x_k + y_k\,\tau_0
\]

\Needspace{5\baselineskip}

\hypertarget{polynomial-drift-removal}{%
\subsection{Polynomial drift removal}\label{polynomial-drift-removal}}

Least squares on a time axis mapped to \([-1, 1]\); the residual
\(\varepsilon(t)\) feeds every estimator:

\[
x(t) = a_0 + a_1 t + a_2 t^2 + \varepsilon(t), \qquad
D = \frac{dy}{dt} = 2 a_2 \quad [\mathrm{s^{-1}}]
\]

\Needspace{5\baselineskip}

\hypertarget{linear-frequency-drift-in-the-allan-deviation}{%
\subsection{Linear frequency drift in the Allan deviation}\label{linear-frequency-drift-in-the-allan-deviation}}

For \(y(t) = D t\):

\[
\sigma_y(\tau) = \frac{D\,\tau}{\sqrt{2}}, \qquad
\sigma_H(\tau) = 0
\]

\Needspace{8\baselineskip}

\hypertarget{validation-performed-003}{%
\section{Validation performed (0.0.3)}\label{validation-performed-003}}

Run with \texttt{run\_\allowbreak{}all.py} on Python 3.8.20 (NumPy 1.24.4, SciPy 1.10.1, PySide6
6.6.3.1), 3.12.13 and 3.13.16 (NumPy 2.5, PySide6 6.12), Qt.py 2.0.7; Veusz
3.6.2 (Ubuntu) and 4.2.1 (PyPI source build).

\begin{longtable}[]{@{}
  >{\raggedright\arraybackslash}p{(\columnwidth - 4\tabcolsep) * \real{0.1451}}
  >{\raggedright\arraybackslash}p{(\columnwidth - 4\tabcolsep) * \real{0.7689}}
  >{\raggedright\arraybackslash}p{(\columnwidth - 4\tabcolsep) * \real{0.0860}}@{}}
\toprule\noalign{}
\begin{minipage}[b]{\linewidth}\raggedright
Script
\end{minipage} & \begin{minipage}[b]{\linewidth}\raggedright
Checks
\end{minipage} & \begin{minipage}[b]{\linewidth}\raggedright
Result
\end{minipage} \\
\midrule\noalign{}
\endhead
\bottomrule\noalign{}
\endlastfoot
\texttt{validate\_\allowbreak{}core.py} & 11 tests, 35 checks: OADEV/MDEV/OHDEV/TDEV = allantools (0.000000 \%); Thêo1 = \texttt{allantools.\allowbreak{}theo1} (≤ 2.5 × 10⁻¹⁶); Thêo1/OADEV(τ) = 1.000\ldots0.971; ThêoH grid per Howe, 75.0 \% coverage; parallel = serial; BLAS-free kernel; frequency CSVs; drift detection; drift fit & PASS \\
\texttt{validate\_\allowbreak{}phase\_\allowbreak{}noise.py} & Slopes −19.80/−29.88/−40.06 dB/decade; noise ID 283/285 = 99 \%; S\_y formula & PASS \\
\texttt{validate\_\allowbreak{}sa\_\allowbreak{}convert.py} & 10 checks incl. decimal-comma R\&S file, bit-exact 10 GHz round trip, 12-file parallel batch & PASS \\
\texttt{validate\_\allowbreak{}stable32\_\allowbreak{}export.py} & 16 checks: read-back, gaps, no partial/stale files, table & PASS \\
\texttt{validate\_\allowbreak{}veusz\_\allowbreak{}export.py} & 26 checks incl. 300 000-point envelope and \texttt{slope\_\allowbreak{}hdev} & PASS \\
\texttt{validate\_\allowbreak{}vszh5\_\allowbreak{}in\_\allowbreak{}veusz.py} & 22/22 with Veusz 3.6.2 and 4.2.1 (loader, bit-exact, tags, references, renders; 250 000-point document) & PASS \\
\texttt{validate\_\allowbreak{}gui.py} & 26 checks: menus, light/dark, worker thread, stale discard, reload, \texttt{freq\_\allowbreak{}hz}, all exports, render, decimal-comma trace, Close Data & PASS \\
Spyder cells & Every \texttt{.py} file run cell by cell under Spyder\textquotesingle s \texttt{runcell} rules & PASS \\
\end{longtable}

\textbf{Not verified:} Windows (paths, the Veusz installer in \texttt{--listen} mode,
spawn pools from a frozen app), a real \texttt{mamba env create}, the GPU path,
instrument data, Stable32 reading the files, machines with more than 2 CPUs.

\Needspace{8\baselineskip}

\hypertarget{notes-for-the-shared-writer-vszh5_writerpy-005}{%
\section{\texorpdfstring{Notes for the shared writer \texttt{vszh5\_writer.py} 0.0.5}{Notes for the shared writer vszh5\_writer.py 0.0.5}}\label{notes-for-the-shared-writer-vszh5_writerpy-005}}

Found while testing this package; the shared module is unchanged here and
\texttt{veusz\_\allowbreak{}export.\allowbreak{}render\_\allowbreak{}vszh5()} works around all four. Worth fixing in the next
writer release (QuickSyn):

\begin{enumerate}
\def\labelenumi{\arabic{enumi}.}
\tightlist
\item
  \texttt{render\_\allowbreak{}pages(png\_\allowbreak{}pages=\allowbreak{}"first")} writes \textbf{no PNG with Veusz 4.2.1}: the
  listen script limits the loop with \texttt{min(1,\allowbreak{} len(names))}, and Veusz 4.2.1
  evaluates scripts with NumPy\textquotesingle s names imported, so \texttt{min} is
  \texttt{numpy.\allowbreak{}min(1,\allowbreak{} axis=\allowbreak{}n)} → \texttt{AxisError}; \texttt{--listen} prints the error and
  still prints the done marker, so the call reports success.
  \texttt{(1 if names else 0)} avoids it. QuickSyn\textquotesingle s export default is
  \texttt{--png-pages first}.
\item
  The PDF is named \texttt{stem.\allowbreak{}with\_\allowbreak{}suffix(".\allowbreak{}pdf")}, so \texttt{run v0.\allowbreak{}0.\allowbreak{}2.\allowbreak{}vszh5} gives
  \texttt{run v0.\allowbreak{}0.pdf}; use \texttt{stem.\allowbreak{}parent /\allowbreak{} (stem.\allowbreak{}name + ".\allowbreak{}pdf")}.
\item
  The returned list holds only files that did not exist before the call, so
  a re-render returns \texttt{[]}.
\item
  \texttt{find\_\allowbreak{}veusz\_\allowbreak{}exe()}/messages hard-code \texttt{QUICKSYN\_\allowbreak{}VEUSZ}; \texttt{render\_\allowbreak{}pages()}
  has no \texttt{env\_\allowbreak{}vars} parameter. Minor: \texttt{local\_\allowbreak{}utc\_\allowbreak{}offsets()} assumes offsets
  change on whole UTC hours, not true for half-hour zones.
\end{enumerate}

\Needspace{8\baselineskip}

\hypertarget{open-items}{%
\section{Open items}\label{open-items}}

\begin{enumerate}
\def\labelenumi{\arabic{enumi}.}
\tightlist
\item
  Confirm GPL-3.0-or-later for this package (needed because it ships the
  GPL \texttt{vszh5\_\allowbreak{}writer.py}) and NRAO/GBO release approval.
\item
  Stable32 time-tag convention (data gap above); test the files in Stable32.
\item
  Windows test: environment, GUI, render, process pools.
\item
  A GPU run of the CuPy path.
\item
  Keysight and R\&S export samples from your instruments for the converter.
\end{enumerate}

\Needspace{8\baselineskip}

\hypertarget{ieee-reference-table}{%
\section{IEEE reference table}\label{ieee-reference-table}}

"Checked" = how the entry was verified in this revision.

\begin{longtable}[]{@{}
  >{\raggedright\arraybackslash}p{(\columnwidth - 6\tabcolsep) * \real{0.0640}}
  >{\raggedright\arraybackslash}p{(\columnwidth - 6\tabcolsep) * \real{0.4523}}
  >{\raggedright\arraybackslash}p{(\columnwidth - 6\tabcolsep) * \real{0.3060}}
  >{\raggedright\arraybackslash}p{(\columnwidth - 6\tabcolsep) * \real{0.1777}}@{}}
\toprule\noalign{}
\begin{minipage}[b]{\linewidth}\raggedright
\#
\end{minipage} & \begin{minipage}[b]{\linewidth}\raggedright
Reference
\end{minipage} & \begin{minipage}[b]{\linewidth}\raggedright
URL
\end{minipage} & \begin{minipage}[b]{\linewidth}\raggedright
Checked
\end{minipage} \\
\midrule\noalign{}
\endhead
\bottomrule\noalign{}
\endlastfoot
{[}1{]} & W. J. Riley, "Handbook of Frequency Stability Analysis," NIST Special Publication 1065, Jul. 2008. & \url{https://tf.nist.gov/general/pdf/2220.pdf} & Document read (sections, equation numbers) \\
{[}2{]} & G. Kamas and M. A. Lombardi, "Time and Frequency Users Manual," NIST Special Publication 559 (Rev. 1990), Sep. 1990. & \url{https://tf.nist.gov/general/pdf/461.pdf} & First page read. 0.0.1 merged it with {[}10{]} \\
{[}3{]} & \emph{IEEE Standard Definitions of Physical Quantities for Fundamental Frequency and Time Metrology --- Random Instabilities}, IEEE Std 1139-2008, 2009, doi: 10.1109/IEEESTD.2008.4797525. & \url{https://ieeexplore.ieee.org/document/4797525} & DOI matches the Xplore record \\
{[}4{]} & D. A. Howe, "ThêoH: a hybrid, high-confidence statistic that improves on the Allan deviation," \emph{Metrologia}, vol. 43, no. 4, pp. S322--S331, 2006, doi: 10.1088/0026-1394/43/4/S20. & \url{https://doi.org/10.1088/0026-1394/43/4/S20} & Crossref record and paper. 0.0.1: DOI \texttt{…S17}, pp. S322--S327 \\
{[}5{]} & D. A. Howe and T. N. Tasset, "Theo1: characterization of very long-term frequency stability," in \emph{Proc. 18th European Frequency and Time Forum}, 2004, pp. 725--728. & \url{https://tf.nist.gov/general/pdf/1990.pdf} & NIST record and paper (τ = 0.75 m τ₀). 0.0.1 named "F. Vernotte" \\
{[}6{]} & D. A. Howe, J. McGee-Taylor and T. Tasset, "ThêoH bias-removal method," in \emph{Proc. 2006 IEEE Int. Frequency Control Symp.}, Miami, FL, USA, Jun. 2006, pp. 788--792. & \url{https://www.nist.gov/publications/theoh-bias-removal-method} & NIST record. 0.0.1 cited "Purdue OxideMEMS Lab" \\
{[}7{]} & C. A. Greenhall and W. J. Riley, "Uncertainty of stability variances based on finite differences," in \emph{Proc. 35th Annu. Precise Time and Time Interval (PTTI) Meeting}, San Diego, CA, USA, Dec. 2003. & \url{https://ntrs.nasa.gov/citations/20050061319} & NASA NTRS record; allantools refs. 0.0.1: Riley first, "36th" \\
{[}8{]} & D. W. Allan, "Statistics of atomic frequency standards," \emph{Proc. IEEE}, vol. 54, no. 2, pp. 221--230, Feb. 1966, doi: 10.1109/PROC.1966.4634. & \url{https://doi.org/10.1109/PROC.1966.4634} & As received (standard citation) \\
{[}9{]} & J. A. Barnes \emph{et al.}, "Characterization of frequency stability," \emph{IEEE Trans. Instrum. Meas.}, vol. IM-20, no. 2, pp. 105--120, May 1971; also NBS Technical Note 394, Oct. 1970. & \url{https://nvlpubs.nist.gov/nistpubs/Legacy/TN/nbstechnicalnote394.pdf} & TN 394 location found; IEEE DOI not checked \\
{[}10{]} & D. B. Sullivan, D. W. Allan, D. A. Howe and F. L. Walls, Eds., "Characterization of Clocks and Oscillators," NIST Technical Note 1337, Mar. 1990. & \url{https://nvlpubs.nist.gov/nistpubs/Legacy/TN/nbstechnicalnote1337.pdf} & Location found in a search; title the 0.0.1 README attached to SP 559 \\
{[}11{]} & P. Welch, "The use of fast Fourier transform for the estimation of power spectra: a method based on time averaging over short, modified periodograms," \emph{IEEE Trans. Audio Electroacoust.}, vol. 15, no. 2, pp. 70--73, Jun. 1967, doi: 10.1109/TAU.1967.1161901. & \url{https://doi.org/10.1109/TAU.1967.1161901} & As received; title completed \\
{[}12{]} & E. Rubiola, \emph{Phase Noise and Frequency Stability in Oscillators}. Cambridge, U.K.: Cambridge Univ. Press, 2008, doi: 10.1017/CBO9780511812798. & \url{https://doi.org/10.1017/CBO9780511812798} & As received \\
{[}13{]} & A. Wallin \emph{et al.}, "allantools," documentation and source (commit ddc5bb5, May 2026). & \url{https://github.com/aewallin/allantools} & Source read; used as reference implementation \\
{[}14{]} & W. J. Riley, "Stable32 User Manual" and "Stable32 Tour," Hamilton Technical Services; Stable32 is distributed by the IEEE UFFC Society. & \url{https://ieee-uffc.org/technical-committees/frequency-control/educational-resources/frequency-control-software} ; \url{https://wriley.com/Tour004.htm} & Pages found; the 0.0.1 URL \texttt{stable32.\allowbreak{}com/\allowbreak{}Manual154.pdf} not verified \\
{[}15{]} & Keysight Technologies, "N9068C Phase Noise Measurement Application --- Measurement Guide" (X-Series). & \url{https://www.keysight.com/} & Not verified (document number) \\
{[}16{]} & Rohde \& Schwarz, "R\&S FSWP Phase Noise Analyzer and VCO Tester --- User Manual." & \url{https://www.rohde-schwarz.com/} & Not verified (document number) \\
{[}17{]} & W. W. Wallace, \texttt{vszh5\_\allowbreak{}writer.py} 0.0.5 (from \texttt{isp\_\allowbreak{}speed\_\allowbreak{}monitor.py} 1.3.0, Adaptive Electrodynamics LLC dba AdaptEd; shipped with QuickSyn 0.0.5), 2026. & Supplied by the author (not public) & SHA-256 matches the QuickSyn record \\
{[}18{]} & Veusz Developers, "Veusz --- a scientific plotting package," and Veusz 4.2.1 source (\texttt{document/\allowbreak{}doc.py}, \texttt{document/\allowbreak{}loader.py}, \texttt{document/\allowbreak{}commandinterface.py}, \texttt{windows/\allowbreak{}mainwindow.py}). & \url{https://veusz.github.io/} ; \url{https://pypi.org/project/veusz/4.2.1/} & Source read; 3.6.2 and 4.2.1 run \\
{[}19{]} & A. Collette \emph{et al.}, "h5py." & \url{https://pypi.org/project/h5py/} & Package used (3.10, 3.11, 3.16) \\
{[}20{]} & M. Ottosson \emph{et al.}, "Qt.py --- minimal Python 2 \& 3 shim around all Qt bindings." & \url{https://github.com/mottosso/Qt.py} & Package used (2.0.7) \\
{[}21{]} & The Qt Company, "Qt for Python (PySide6) documentation." & \url{https://doc.qt.io/qtforpython-6/} & Package used (6.6.3.1, 6.12) \\
{[}22{]} & J. D. Hunter, "Matplotlib: a 2D graphics environment," \emph{Comput. Sci. Eng.}, vol. 9, no. 3, pp. 90--95, 2007, doi: 10.1109/MCSE.2007.55. & \url{https://doi.org/10.1109/MCSE.2007.55} & As received \\
{[}23{]} & P. Virtanen \emph{et al.}, "SciPy 1.0: fundamental algorithms for scientific computing in Python," \emph{Nature Methods}, vol. 17, pp. 261--272, 2020, doi: 10.1038/s41592-019-0686-2. & \url{https://doi.org/10.1038/s41592-019-0686-2} & Standard citation \\
{[}24{]} & C. R. Harris \emph{et al.}, "Array programming with NumPy," \emph{Nature}, vol. 585, pp. 357--362, 2020, doi: 10.1038/s41586-020-2649-2. & \url{https://doi.org/10.1038/s41586-020-2649-2} & Standard citation \\
{[}25{]} & R. Okuta, Y. Unno, D. Nishino, S. Hido and C. Loomis, "CuPy: a NumPy-compatible library for NVIDIA GPU calculations," in \emph{Proc. Workshop on ML Systems (LearningSys)}, NIPS 2017. & \url{https://learningsys.org/nips17/assets/papers/paper_16.pdf} & Standard citation \\
{[}26{]} & joblib developers, "threadpoolctl." & \url{https://github.com/joblib/threadpoolctl} & Package used \\
{[}27{]} & T. Preston-Werner, "Semantic Versioning 2.0.0." & \url{https://semver.org/spec/v2.0.0.html} & --- \\
{[}28{]} & Python Software Foundation, "multiprocessing --- The spawn and forkserver start methods"; Spyder developers, \texttt{spyder-kernels} 3.1.6 source (\texttt{customize/\allowbreak{}namespace\_\allowbreak{}manager.py}). & \url{https://docs.python.org/3/library/multiprocessing.html} ; \url{https://github.com/spyder-ide/spyder-kernels} & Source read \\
\end{longtable}

Removed in 0.0.3: "MHM-2010 / MSM-2010 Hydrogen Maser User Guide, Rev. J" ---
not cited anywhere in the package, and the revision could not be verified.


\end{document}
